\documentclass[12pt,a4paper]{article}

% ------------------------------------------------------------
% Кодировка, язык и шрифты
% ------------------------------------------------------------
\usepackage[T2A]{fontenc}
\usepackage[utf8]{inputenc}
\usepackage[russian]{babel}

% ------------------------------------------------------------
% Математика
% ------------------------------------------------------------
\usepackage{amsmath,amssymb,amsthm,mathtools,bm}

% ------------------------------------------------------------
% Таблицы, рисунки, списки
% ------------------------------------------------------------
\usepackage{graphicx}
\usepackage{booktabs}
\usepackage{tabularx}
\usepackage{longtable}
\usepackage{array}
\usepackage{enumitem}
\usepackage{xcolor}
\usepackage[expansion=false]{microtype}

% ------------------------------------------------------------
% Алгоритмы
% ------------------------------------------------------------
\usepackage{algorithm}
\usepackage{algpseudocode}
\floatname{algorithm}{Алгоритм}

\algrenewcommand\algorithmicrequire{\textbf{Вход:}}
\algrenewcommand\algorithmicensure{\textbf{Выход:}}
\algrenewcommand\algorithmicwhile{\textbf{пока}}
\algrenewcommand\algorithmicdo{\textbf{выполнять}}
\algrenewcommand\algorithmicfor{\textbf{для}}
\algrenewcommand\algorithmicforall{\textbf{для всех}}
\algrenewcommand\algorithmicif{\textbf{если}}
\algrenewcommand\algorithmicthen{\textbf{то}}
\algrenewcommand\algorithmicelse{\textbf{иначе}}
\algrenewcommand\algorithmicreturn{\textbf{вернуть}}

% ------------------------------------------------------------
% Диаграммы
% ------------------------------------------------------------
\usepackage{tikz}
\usetikzlibrary{
    arrows.meta,
    positioning,
    shapes.geometric,
    fit,
    backgrounds
}

% ------------------------------------------------------------
% Геометрия страницы
% ------------------------------------------------------------
\usepackage[
    left=25mm,
    right=20mm,
    top=20mm,
    bottom=22mm
]{geometry}

% ------------------------------------------------------------
% Гиперссылки внутри документа
% ------------------------------------------------------------
\usepackage{hyperref}
\hypersetup{hidelinks}

% ------------------------------------------------------------
% Команды и окружения
% ------------------------------------------------------------
\DeclareMathOperator{\softmax}{softmax}
\DeclareMathOperator{\argmax}{arg\,max}
\DeclareMathOperator{\argmin}{arg\,min}
\DeclareMathOperator{\SiLU}{SiLU}
\DeclareMathOperator{\KL}{KL}
\DeclareMathOperator{\diag}{diag}

\newcommand{\R}{\mathbb{R}}
\newcommand{\V}{\mathcal{V}}
\newcommand{\D}{\mathcal{D}}
\newcommand{\E}{\mathbb{E}}
\newcommand{\1}{\mathbf{1}}

\newtheorem{theorem}{Теорема}[section]
\newtheorem{proposition}[theorem]{Утверждение}
\newtheorem{lemma}[theorem]{Лемма}
\newtheorem{corollary}[theorem]{Следствие}

\theoremstyle{definition}
\newtheorem{definition}[theorem]{Определение}
\newtheorem{example}[theorem]{Пример}

\theoremstyle{remark}
\newtheorem{remark}[theorem]{Замечание}

\setlist[itemize]{topsep=3pt,itemsep=2pt}
\setlist[enumerate]{topsep=3pt,itemsep=2pt}

\tikzset{
    process/.style={
        draw,
        rounded corners,
        align=center,
        minimum height=10mm,
        text width=28mm,
        fill=blue!5
    },
    data/.style={
        draw,
        cylinder,
        shape border rotate=90,
        aspect=0.25,
        align=center,
        minimum height=12mm,
        text width=24mm,
        fill=green!7
    },
    arrow/.style={
        -{Latex[length=2.5mm]},
        thick
    }
}

\title{Билет №100 \\ \small Большие языковые модели. Принципы построения и использования. (ИТМО)}
\author{Сериков Владислав Александрович}
\date{19 июля 2026}

\begin{document}

\maketitle
\tableofcontents
\newpage

% ============================================================
\section{Основные понятия}
% ============================================================

\subsection{Языковая модель}

\begin{definition}
\emph{Языковая модель} --- вероятностная модель, задающая распределение
вероятностей на последовательностях языковых единиц либо условное
распределение следующей единицы при известном контексте.
\end{definition}

Пусть $\V$ --- конечный словарь токенов, а
\[
x_{1:n}=(x_1,x_2,\ldots,x_n), \qquad x_i\in\V,
\]
--- токенизированная последовательность. Авторегрессионная языковая модель
задаёт условные вероятности
\[
p_\theta(x_t\mid x_{<t}),
\qquad
x_{<t}=(x_1,\ldots,x_{t-1}),
\]
где $\theta$ обозначает параметры модели.

Термин \emph{большая языковая модель} (Large Language Model, LLM) не имеет
единого строгого порога по числу параметров. Обычно под LLM понимают
нейросетевую языковую модель, которая:

\begin{itemize}
    \item обучена на большом и разнообразном корпусе;
    \item имеет достаточно большую ёмкость;
    \item переносит знания между задачами;
    \item может решать новые задачи по инструкции или нескольким примерам;
    \item требует существенных вычислительных ресурсов для полного обучения.
\end{itemize}

Размер модели сам по себе недостаточен для описания её качества. Существенны
также объём и качество данных, вычислительный бюджет, архитектура,
оптимизация и последующее обучение.

\subsection{Базовая, инструктивная и диалоговая модели}

Различают несколько стадий одной модели.

\begin{description}
    \item[Базовая модель.] Обучена главным образом предсказывать следующий
    токен. Она умеет продолжать текст, но не обязана точно выполнять
    пользовательские инструкции.

    \item[Инструктивная модель.] Дополнительно обучена на парах
    <<инструкция--ответ>>. Она лучше интерпретирует запрос как задачу.

    \item[Диалоговая модель.] Инструктивная модель, адаптированная к
    многоходовому диалогу, ролям сообщений, ограничениям безопасности
    и сохранению контекста беседы.

    \item[Предметно-адаптированная модель.] Дополнительно обучена на данных
    конкретной области: программировании, медицине, праве, математике,
    технической документации и т.\,п.
\end{description}

\subsection{Параметрические и внешние знания}

Знания LLM могут поступать из двух источников:

\begin{enumerate}
    \item \textbf{параметрическая память} --- закономерности, сохранённые
    в весах нейронной сети во время обучения;
    \item \textbf{непараметрическая память} --- документы, базы данных,
    поисковые индексы и программные инструменты, доступные во время запроса.
\end{enumerate}

Параметрическая память компактна и быстро используется, но её трудно
точечно обновлять. Внешнюю базу знаний проще обновить и проверить, однако
для неё требуются поиск, управление контекстом и контроль качества
извлечённых данных.

\subsection{Жизненный цикл языковой модели}

\begin{figure}[h]
\centering
\resizebox{\textwidth}{!}{
\begin{tikzpicture}[node distance=9mm]
    \node[data] (raw) {Исходные\\данные};
    \node[process, right=of raw] (prep) {Очистка,\\фильтрация,\\дедупликация};
    \node[process, right=of prep] (tok) {Обучение\\токенизатора};
    \node[process, right=of tok] (pre) {Предварительное\\обучение};
    \node[process, right=of pre] (post) {Инструктивное\\и предпочтительное\\обучение};
    \node[process, right=of post] (serve) {Инференс, RAG,\\инструменты};
    \node[process, right=of serve] (eval) {Оценивание\\и мониторинг};

    \draw[arrow] (raw) -- (prep);
    \draw[arrow] (prep) -- (tok);
    \draw[arrow] (tok) -- (pre);
    \draw[arrow] (pre) -- (post);
    \draw[arrow] (post) -- (serve);
    \draw[arrow] (serve) -- (eval);
    \draw[arrow, bend left=35] (eval.north) to
        node[above, align=center]{обратная связь} (prep.north);
\end{tikzpicture}
}
\caption{Типичный жизненный цикл большой языковой модели.}
\end{figure}

% ============================================================
\section{Вероятностные основы языкового моделирования}
% ============================================================

\subsection{Цепное разложение вероятности}

\begin{theorem}[Цепное правило вероятностей]
Пусть $X_1,\ldots,X_n$ --- дискретные случайные величины. Для значений
$x_1,\ldots,x_n$, для которых соответствующие условные вероятности
определены,
\[
\Pr(X_1=x_1,\ldots,X_n=x_n)
=
\Pr(X_1=x_1)
\prod_{t=2}^{n}
\Pr(X_t=x_t\mid X_1=x_1,\ldots,X_{t-1}=x_{t-1}).
\]
\end{theorem}

\begin{proof}
Для двух случайных величин по определению условной вероятности
\[
\Pr(X_1=x_1,X_2=x_2)
=
\Pr(X_1=x_1)
\Pr(X_2=x_2\mid X_1=x_1).
\]

Предположим, что формула доказана для $n-1$ случайной величины. Тогда
\begin{align*}
&\Pr(X_1=x_1,\ldots,X_n=x_n)
\\
&\quad =
\Pr(X_1=x_1,\ldots,X_{n-1}=x_{n-1})
\\
&\qquad\quad{}\times
\Pr(X_n=x_n\mid X_1=x_1,\ldots,X_{n-1}=x_{n-1}).
\end{align*}
Применяя предположение индукции к первому множителю, получаем
\[
\Pr(X_1=x_1)
\prod_{t=2}^{n-1}
\Pr(X_t=x_t\mid X_1=x_1,\ldots,X_{t-1}=x_{t-1})
\]
и умножаем это выражение на условную вероятность для $X_n$.
Следовательно, формула верна для $n$. По математической индукции она
верна для любого $n$.
\end{proof}

Поэтому вероятность последовательности токенов можно записать как
\[
p_\theta(x_{1:n})
=
\prod_{t=1}^{n}p_\theta(x_t\mid x_{<t}),
\]
где при $t=1$ в контекст может входить специальный токен начала
последовательности $\texttt{<BOS>}$.

Это разложение превращает моделирование распределения над всеми текстами в
повторяющуюся задачу классификации: для каждого контекста требуется
предсказать распределение по словарю $\V$.

\subsection{Максимизация правдоподобия}

Пусть обучающий корпус состоит из последовательностей
\[
\D=\{x^{(1)},x^{(2)},\ldots,x^{(M)}\}.
\]
Метод максимального правдоподобия выбирает
\[
\theta^\ast
=
\argmax_\theta
\prod_{m=1}^{M}p_\theta(x^{(m)}).
\]
Логарифм монотонно возрастает, поэтому эквивалентно максимизировать
\[
\sum_{m=1}^{M}\log p_\theta(x^{(m)})
=
\sum_{m=1}^{M}\sum_{t}
\log p_\theta\!\left(x_t^{(m)}\mid x_{<t}^{(m)}\right).
\]

На практике минимизируют отрицательное логарифмическое правдоподобие:
\[
\mathcal{L}_{\mathrm{NLL}}(\theta)
=
-\frac{1}{N}
\sum_{m=1}^{M}\sum_t
\log p_\theta\!\left(x_t^{(m)}\mid x_{<t}^{(m)}\right),
\]
где $N$ --- число целевых токенов.

При обучении используется \emph{teacher forcing}: для предсказания токена
$x_t$ модели передают истинный префикс $x_{<t}$, а не ранее сгенерированные
ею токены.

Например, для последовательности
\[
[\texttt{<BOS>},\ \texttt{я},\ \texttt{люблю},\
 \texttt{чай},\ \texttt{<EOS>}]
\]
входы и цели сдвигаются на одну позицию:
\[
\begin{array}{c|cccc}
\text{вход} &
\texttt{<BOS>} & \texttt{я} & \texttt{люблю} & \texttt{чай}
\\
\hline
\text{цель} &
\texttt{я} & \texttt{люблю} & \texttt{чай} & \texttt{<EOS>}.
\end{array}
\]

\subsection{Энтропия, кросс-энтропия и дивергенция Кульбака--Лейблера}

Пусть $P=(p_1,\ldots,p_K)$ и $Q=(q_1,\ldots,q_K)$ --- распределения на
конечном множестве, причём $p_i>0$ и $q_i>0$.

Энтропия распределения $P$:
\[
H(P)=-\sum_{i=1}^{K}p_i\log p_i.
\]

Кросс-энтропия:
\[
H(P,Q)=-\sum_{i=1}^{K}p_i\log q_i.
\]

Дивергенция Кульбака--Лейблера:
\[
\KL(P\Vert Q)
=
\sum_{i=1}^{K}p_i\log\frac{p_i}{q_i}.
\]

Между ними выполняется тождество
\[
H(P,Q)=H(P)+\KL(P\Vert Q).
\]

\begin{theorem}[Неравенство Гиббса]
Для любых строго положительных конечных распределений $P$ и $Q$
\[
\KL(P\Vert Q)\geq 0.
\]
Равенство выполняется тогда и только тогда, когда $P=Q$.
\end{theorem}

\begin{proof}
Для любого $u>0$ справедливо
\[
\log u\leq u-1,
\]
причём равенство достигается только при $u=1$. Подставим
$u=q_i/p_i$:
\[
-\log\frac{q_i}{p_i}
\geq
1-\frac{q_i}{p_i}.
\]
Умножим на $p_i>0$:
\[
p_i\log\frac{p_i}{q_i}\geq p_i-q_i.
\]
Суммируя по $i$, получаем
\[
\KL(P\Vert Q)
=
\sum_i p_i\log\frac{p_i}{q_i}
\geq
\sum_i(p_i-q_i)=1-1=0.
\]

Равенство возможно только тогда, когда для каждого $i$
\[
\frac{q_i}{p_i}=1,
\]
то есть $q_i=p_i$. Следовательно, равенство достигается тогда и только
тогда, когда $P=Q$.
\end{proof}

\begin{corollary}
Кросс-энтропия $H(P,Q)$ как функция $Q$ минимальна при $Q=P$.
\end{corollary}

\begin{proof}
Так как
\[
H(P,Q)=H(P)+\KL(P\Vert Q),
\]
а $H(P)$ не зависит от $Q$, минимум достигается при минимуме
$\KL(P\Vert Q)$. По неравенству Гиббса этот минимум равен нулю и
достигается только при $Q=P$.
\end{proof}

Следовательно, минимизация кросс-энтропии заставляет модель приближать
истинное условное распределение следующего токена.

\begin{example}
После контекста \texttt{я люблю} в обучающей выборке наблюдались:
\[
\texttt{чай}:3\text{ раза},
\qquad
\texttt{кофе}:1\text{ раз}.
\]
Эмпирическое распределение равно
\[
\widehat{P}(\texttt{чай})=\frac{3}{4},
\qquad
\widehat{P}(\texttt{кофе})=\frac{1}{4}.
\]
Если контекст рассматривать независимо от других контекстов, максимум
правдоподобия достигается при таком же предсказываемом распределении.
Нейронная модель дополнительно разделяет параметры между множеством
контекстов и благодаря этому может обобщать закономерности.
\end{example}

\subsection{Перплексия}

Средней токенной кросс-энтропии соответствует метрика
\emph{перплексия}:
\[
\operatorname{PPL}
=
\exp\left(
-\frac{1}{N}\sum_{t=1}^{N}
\log p_\theta(x_t\mid x_{<t})
\right).
\]

Чем меньше перплексия, тем большую вероятность модель в среднем
приписывает правильным токенам. Условно её можно интерпретировать как
эффективное число равновероятных продолжений на каждом шаге.

Перплексия имеет ограничения:

\begin{itemize}
    \item зависит от токенизатора;
    \item не измеряет напрямую фактическую корректность;
    \item не гарантирует выполнение инструкций;
    \item не полностью характеризует качество длинного текста;
    \item может улучшаться из-за утечки тестовых данных.
\end{itemize}

% ============================================================
\section{Токенизация}
% ============================================================

\subsection{Назначение токенизатора}

Нейронная сеть работает с числами, поэтому текст преобразуется в
последовательность идентификаторов:
\[
T:\text{строка}\longrightarrow(x_1,\ldots,x_n),
\qquad x_i\in\{0,\ldots,|\V|-1\}.
\]

Обратное преобразование называется детокенизацией:
\[
T^{-1}:(x_1,\ldots,x_n)\longrightarrow\text{строка}.
\]

\begin{center}
\begin{tabularx}{\textwidth}{>{\bfseries}l X X}
\toprule
Единица & Преимущества & Недостатки \\
\midrule
Слово
& Короткие последовательности, естественная интерпретация
& Огромный словарь, проблема неизвестных слов и словоформ
\\
Символ
& Малый словарь, отсутствие неизвестных слов
& Очень длинные последовательности
\\
Подслово
& Компромисс между словами и символами
& Разбиение не всегда совпадает с морфемами
\\
Байт
& Можно представить любую строку
& Некоторые языки и символы кодируются длиннее
\\
\bottomrule
\end{tabularx}
\end{center}

В современных LLM обычно применяют подсловные или байтовые токены.

\subsection{Byte Pair Encoding}

Изначально Byte Pair Encoding был алгоритмом сжатия, но в языковом
моделировании его используют для построения подсловного словаря.

\begin{algorithm}[h]
\caption{Обучение подсловного BPE}
\begin{algorithmic}[1]
\Require словарь слов с частотами $C$; число объединений $M$
\Ensure упорядоченный список правил объединения $R$
\State представить каждое слово последовательностью начальных единиц
\State $R\gets [\ ]$
\For{$m=1,\ldots,M$}
    \State подсчитать частоты всех соседних пар с учётом частот слов
    \State выбрать наиболее частую пару $(a,b)$
    \State добавить правило $(a,b)\mapsto ab$ в конец $R$
    \State заменить все соседние вхождения $(a,b)$ единицей $ab$
\EndFor
\State \Return $R$
\end{algorithmic}
\end{algorithm}

При токенизации нового текста начальные единицы последовательно
объединяются в порядке, заданном списком правил $R$.

\begin{example}
Рассмотрим корпус:
\[
\texttt{кот}\times 5,\qquad
\texttt{коты}\times 3,\qquad
\texttt{котик}\times 2.
\]
Начальное символьное представление:
\[
\begin{aligned}
&\texttt{к о т </w>} &&\text{с частотой }5,\\
&\texttt{к о т ы </w>} &&\text{с частотой }3,\\
&\texttt{к о т и к </w>} &&\text{с частотой }2.
\end{aligned}
\]

Пары $\texttt{к о}$ и $\texttt{о т}$ встречаются по $10$ раз.
Пусть при равенстве выбрана пара $\texttt{к о}$. Первое правило:
\[
(\texttt{к},\texttt{о})\mapsto\texttt{ко}.
\]
После замены пара $\texttt{ко т}$ встречается $10$ раз. Второе правило:
\[
(\texttt{ко},\texttt{т})\mapsto\texttt{кот}.
\]

После двух объединений слова представляются как
\[
\texttt{кот},\qquad
\texttt{кот ы},\qquad
\texttt{кот и к}.
\]
Частая основа получила собственный токен, а окончания и редкие части
остались более мелкими единицами.
\end{example}

\subsection{Другие подходы}

\begin{description}
    \item[WordPiece.] Выбирает объединения с учётом статистического
    критерия, связанного с вероятностью единиц, а не только с абсолютной
    частотой пары.

    \item[Unigram Language Model.] Начинает с большого набора возможных
    подслов и итеративно удаляет наименее полезные единицы, оптимизируя
    правдоподобие сегментаций.

    \item[SentencePiece.] Обучается непосредственно на необработанных
    строках и рассматривает пробел как обычный специальный символ.
    Это уменьшает зависимость от языково-специфичного предварительного
    разбиения на слова.

    \item[Byte-level BPE.] Сначала кодирует текст байтами, а затем обучает
    BPE. Любая входная строка представима без специального неизвестного
    токена.
\end{description}

\subsection{Специальные токены}

Типичные специальные токены:

\begin{itemize}
    \item $\texttt{<BOS>}$ --- начало последовательности;
    \item $\texttt{<EOS>}$ --- конец последовательности;
    \item $\texttt{<PAD>}$ --- дополнение последовательностей в пакете;
    \item $\texttt{<UNK>}$ --- неизвестная единица, если она допускается;
    \item токены ролей: система, пользователь, ассистент;
    \item токены начала и конца документа;
    \item токены вызова инструментов или структурированного вывода.
\end{itemize}

Токены ролей и разделителей должны использоваться точно в том формате,
который применялся при последующем обучении модели. Несоответствие
диалогового шаблона может существенно ухудшить поведение.

\subsection{Выбор размера словаря}

Большой словарь:

\begin{itemize}
    \item уменьшает среднюю длину последовательности;
    \item увеличивает матрицу эмбеддингов и выходной слой;
    \item может содержать много редких единиц.
\end{itemize}

Малый словарь:

\begin{itemize}
    \item уменьшает размер эмбеддингов;
    \item увеличивает длину контекста в токенах;
    \item повышает стоимость механизма внимания;
    \item может ухудшать обработку отдельных языков и терминов.
\end{itemize}

Токенизатор является частью модели. Изменение словаря после обучения
требует изменения матрицы эмбеддингов и обычно дополнительного обучения.

% ============================================================
\section{Архитектура Transformer}
% ============================================================

\subsection{Эмбеддинги токенов}

Пусть $V=|\V|$ --- размер словаря, $d$ --- размерность скрытого
представления. Матрица эмбеддингов
\[
E\in\R^{V\times d}
\]
сопоставляет токену $x_i$ вектор
\[
e_i=E[x_i]\in\R^d.
\]

Для последовательности длины $n$ получаем матрицу
\[
X=
\begin{bmatrix}
e_1^\top\\
e_2^\top\\
\vdots\\
e_n^\top
\end{bmatrix}
\in\R^{n\times d}.
\]

Сам по себе механизм внимания без позиционной информации не знает
порядка токенов. Поэтому к эмбеддингам добавляется или иным способом
встраивается информация о позиции.

\subsection{Позиционные представления}

\subsubsection{Синусоидальное кодирование}

В исходной архитектуре Transformer использовалось
\[
\operatorname{PE}(m,2j)
=
\sin\left(
\frac{m}{10000^{2j/d}}
\right),
\]
\[
\operatorname{PE}(m,2j+1)
=
\cos\left(
\frac{m}{10000^{2j/d}}
\right),
\]
где $m$ --- номер позиции, $j$ --- номер пары координат.

Вход блока:
\[
h_m^{(0)}=E[x_m]+\operatorname{PE}(m).
\]

\subsubsection{Обучаемые позиционные эмбеддинги}

Для каждой позиции хранится обучаемый вектор:
\[
P\in\R^{n_{\max}\times d},
\qquad
h_m^{(0)}=E[x_m]+P_m.
\]
Недостаток --- жёсткая привязка к максимальной длине, использованной
при обучении.

\subsubsection{Rotary Position Embedding}

В RoPE пары координат запросов и ключей поворачиваются на угол,
зависящий от позиции. Для одной двумерной пары:
\[
R_m=
\begin{bmatrix}
\cos(m\omega)&-\sin(m\omega)\\
\sin(m\omega)& \cos(m\omega)
\end{bmatrix}.
\]
Запрос и ключ преобразуются:
\[
q_m'=R_mq_m,\qquad k_n'=R_nk_n.
\]
Тогда
\[
\langle q_m',k_n'\rangle
=
q_m^\top R_m^\top R_n k_n
=
q_m^\top R_{n-m}k_n.
\]
Таким образом, скалярное произведение зависит от относительного смещения
$n-m$.

\subsection{Scaled Dot-Product Attention}

Пусть
\[
X\in\R^{n\times d}.
\]
Линейными преобразованиями строятся запросы, ключи и значения:
\[
Q=XW_Q,\qquad
K=XW_K,\qquad
V=XW_V,
\]
где
\[
W_Q,W_K\in\R^{d\times d_k},
\qquad
W_V\in\R^{d\times d_v}.
\]

Матрица внимания:
\[
A
=
\softmax_{\mathrm{row}}
\left(
\frac{QK^\top}{\sqrt{d_k}}+M
\right),
\]
а результат:
\[
\operatorname{Attention}(Q,K,V)=AV.
\]

Деление на $\sqrt{d_k}$ ограничивает масштаб скалярных произведений.
Без масштабирования при больших $d_k$ значения могут попадать в область,
где softmax близок к насыщению и его градиенты малы.

Softmax для вектора $z\in\R^m$:
\[
\softmax(z)_i
=
\frac{\exp(z_i)}
{\sum_{j=1}^{m}\exp(z_j)}.
\]

\begin{example}[Численный пример внимания]
Пусть
\[
q=(1,0),\qquad
k_1=(1,0),\qquad
k_2=(0,1),
\]
\[
v_1=(1,2),\qquad
v_2=(3,0),\qquad d_k=2.
\]
Оценки соответствия:
\[
s_1=\frac{q^\top k_1}{\sqrt{2}}=\frac{1}{\sqrt{2}},
\qquad
s_2=\frac{q^\top k_2}{\sqrt{2}}=0.
\]
После softmax:
\[
\alpha_1
=
\frac{e^{1/\sqrt{2}}}{e^{1/\sqrt{2}}+1}
\approx 0.67,
\qquad
\alpha_2\approx 0.33.
\]
Выход:
\[
o=\alpha_1v_1+\alpha_2v_2
\approx
0.67(1,2)+0.33(3,0)
=
(1.66,1.34).
\]
Запрос сильнее соответствует первому ключу, поэтому первый вектор
значений получает больший вес.
\end{example}

\begin{proposition}[Выход внимания как выпуклая комбинация]
Для одной строки обычного softmax-внимания без dropout выход
\[
o_i=\sum_{j=1}^{n}\alpha_{ij}v_j
\]
является выпуклой комбинацией векторов $v_1,\ldots,v_n$.
\end{proposition}

\begin{proof}
По определению softmax
\[
\alpha_{ij}
=
\frac{\exp(s_{ij})}
{\sum_{\ell=1}^{n}\exp(s_{i\ell})}.
\]
Следовательно,
\[
\alpha_{ij}>0
\]
и
\[
\sum_{j=1}^{n}\alpha_{ij}=1.
\]
По определению выпуклой комбинации сумма векторов с неотрицательными
коэффициентами, сумма которых равна единице, является выпуклой
комбинацией. Поэтому $o_i$ принадлежит выпуклой оболочке
$\{v_1,\ldots,v_n\}$.
\end{proof}

\subsection{Каузальная маска}

В авторегрессионной модели позиция $i$ не должна видеть будущие токены.
Используется маска
\[
M_{ij}
=
\begin{cases}
0,&j\leq i,\\
-\infty,&j>i.
\end{cases}
\]
После softmax значения, соответствующие $j>i$, получают нулевую
вероятность.

Для последовательности длины $4$ допустимые связи имеют вид
\[
\begin{array}{c|cccc}
 &1&2&3&4\\
\hline
1&\checkmark&&&\\
2&\checkmark&\checkmark&&\\
3&\checkmark&\checkmark&\checkmark&\\
4&\checkmark&\checkmark&\checkmark&\checkmark
\end{array}
\]

\begin{theorem}[Каузальность decoder-only Transformer]
Рассмотрим Transformer, в котором:

\begin{enumerate}
    \item начальное представление позиции $i$ зависит только от токена
    $x_i$ и номера позиции $i$;
    \item каждый слой самовнимания использует каузальную маску;
    \item нормализация и полносвязная сеть применяются отдельно к каждой
    позиции.
\end{enumerate}

Тогда скрытое представление позиции $i$ на любом слое зависит только от
префикса $x_{1:i}$ и не зависит от $x_{i+1:n}$.
\end{theorem}

\begin{proof}
Проведём индукцию по номеру слоя.

На слое $0$ представление $h_i^{(0)}$ зависит только от $x_i$ и позиции
$i$. Следовательно, оно зависит только от $x_{1:i}$.

Предположим, что для некоторого слоя $\ell$ каждое представление
$h_j^{(\ell)}$ зависит только от $x_{1:j}$.

На следующем слое позиция $i$ в механизме внимания может использовать
только позиции $j\leq i$. По предположению индукции
$h_j^{(\ell)}$ зависит только от $x_{1:j}$, а этот префикс является
частью $x_{1:i}$. Поэтому результат внимания в позиции $i$ зависит
только от $x_{1:i}$.

Остаточная связь объединяет этот результат с $h_i^{(\ell)}$, которое
также зависит только от $x_{1:i}$. Нормализация и полносвязная сеть
применяются независимо к позиции $i$ и не добавляют зависимости от
других токенов. Значит, $h_i^{(\ell+1)}$ зависит только от $x_{1:i}$.

По индукции утверждение верно для всех слоёв.
\end{proof}

\subsection{Multi-Head Attention}

Одна голова внимания может специализироваться только на одном
представлении соответствия. В многоголовом внимании вычисляются
несколько голов:
\[
\operatorname{head}_r
=
\operatorname{Attention}
\left(
XW_Q^{(r)},
XW_K^{(r)},
XW_V^{(r)}
\right),
\]
после чего результаты объединяются:
\[
\operatorname{MHA}(X)
=
\operatorname{Concat}
(\operatorname{head}_1,\ldots,\operatorname{head}_h)W_O.
\]

Обычно
\[
d_k=d_v=\frac{d}{h}.
\]

Варианты:

\begin{description}
    \item[Multi-Head Attention, MHA.] Каждая голова имеет собственные
    ключи и значения.

    \item[Multi-Query Attention, MQA.] Все головы запросов разделяют
    одну голову ключей и значений. Это уменьшает KV-кэш.

    \item[Grouped-Query Attention, GQA.] Головы запросов разбиваются на
    группы; внутри группы ключи и значения разделяются. Это компромисс
    между MHA и MQA.
\end{description}

\subsection{Эквивариантность внимания к перестановкам}

\begin{theorem}
Самовнимание без позиционной информации и без каузальной маски
эквивариантно к перестановке входных токенов. То есть для любой матрицы
перестановки $P$
\[
\operatorname{Attention}(PX)=
P\operatorname{Attention}(X).
\]
\end{theorem}

\begin{proof}
Пусть
\[
Q=XW_Q,\qquad K=XW_K,\qquad V=XW_V.
\]
После перестановки строк входа:
\[
Q'=PXW_Q=PQ,\qquad
K'=PK,\qquad
V'=PV.
\]
Матрица оценок:
\[
Q'K'^\top
=
(PQ)(PK)^\top
=
PQK^\top P^\top.
\]

Одновременная перестановка строк и столбцов только переупорядочивает
элементы каждой строки. Поэтому построчный softmax удовлетворяет
\[
\softmax_{\mathrm{row}}(PAP^\top)
=
P\softmax_{\mathrm{row}}(A)P^\top.
\]
Следовательно,
\begin{align*}
\operatorname{Attention}(PX)
&=
\softmax_{\mathrm{row}}
\left(
\frac{PQK^\top P^\top}{\sqrt{d_k}}
\right)PV
\\
&=
P\softmax_{\mathrm{row}}
\left(
\frac{QK^\top}{\sqrt{d_k}}
\right)P^\top PV
\\
&=
P\softmax_{\mathrm{row}}
\left(
\frac{QK^\top}{\sqrt{d_k}}
\right)V
\\
&=
P\operatorname{Attention}(X).
\end{align*}
Таким образом, выход переставляется так же, как вход.
\end{proof}

Позиционные представления необходимы именно для нарушения этой симметрии
и кодирования порядка.

\subsection{Нормализация и остаточные связи}

Остаточная связь имеет вид
\[
y=x+F(x).
\]
Она создаёт короткий путь для передачи сигнала и градиента через
глубокую сеть.

Layer Normalization:
\[
\mu(x)=\frac{1}{d}\sum_{j=1}^{d}x_j,
\qquad
\sigma^2(x)=\frac{1}{d}\sum_{j=1}^{d}(x_j-\mu(x))^2,
\]
\[
\operatorname{LayerNorm}(x)
=
\gamma\odot
\frac{x-\mu(x)\1}{\sqrt{\sigma^2(x)+\varepsilon}}
+\beta.
\]

RMSNorm не вычитает среднее:
\[
\operatorname{RMS}(x)
=
\sqrt{\frac{1}{d}\sum_{j=1}^{d}x_j^2+\varepsilon},
\]
\[
\operatorname{RMSNorm}(x)
=
\gamma\odot\frac{x}{\operatorname{RMS}(x)}.
\]

В Pre-Norm Transformer нормализация выполняется перед подслоем:
\[
y=x+\operatorname{Attention}(\operatorname{Norm}(x)),
\]
\[
z=y+\operatorname{FFN}(\operatorname{Norm}(y)).
\]
Такая схема обычно облегчает обучение глубоких моделей.

\subsection{Полносвязный подслой}

Стандартная позиционно-независимая сеть:
\[
\operatorname{FFN}(x)
=
W_2\phi(W_1x+b_1)+b_2,
\]
где
\[
W_1\in\R^{d_{\mathrm{ff}}\times d},
\qquad
W_2\in\R^{d\times d_{\mathrm{ff}}}.
\]

Часто
\[
d_{\mathrm{ff}}\approx 4d.
\]

Вариант SwiGLU:
\[
\operatorname{SwiGLU}(x)
=
W_o
\left[
\SiLU(W_gx)
\odot
(W_ux)
\right],
\]
где
\[
\SiLU(z)=z\sigma(z).
\]
Знак $\odot$ обозначает покомпонентное умножение.

Механизм внимания смешивает информацию между позициями, а FFN
нелинейно преобразует признаки внутри каждой позиции.

\subsection{Полный decoder-only блок}

\begin{figure}[h]
\centering
\begin{tikzpicture}[node distance=7mm]
    \node[process] (input) {Скрытые\\состояния $H$};
    \node[process, below=of input] (norm1) {RMSNorm или\\LayerNorm};
    \node[process, below=of norm1] (attn) {Каузальное\\самовнимание};
    \node[process, below=of attn] (add1) {Остаточное\\сложение};
    \node[process, below=of add1] (norm2) {RMSNorm или\\LayerNorm};
    \node[process, below=of norm2] (ffn) {FFN или\\SwiGLU};
    \node[process, below=of ffn] (add2) {Остаточное\\сложение};
    \node[process, below=of add2] (output) {Выход блока};

    \draw[arrow] (input) -- (norm1);
    \draw[arrow] (norm1) -- (attn);
    \draw[arrow] (attn) -- (add1);
    \draw[arrow] (add1) -- (norm2);
    \draw[arrow] (norm2) -- (ffn);
    \draw[arrow] (ffn) -- (add2);
    \draw[arrow] (add2) -- (output);

    \draw[arrow, bend right=55] (input.west) to (add1.west);
    \draw[arrow, bend right=55] (add1.west) to (add2.west);
\end{tikzpicture}
\caption{Структура Pre-Norm decoder-only блока.}
\end{figure}

После $L$ блоков применяется финальная нормализация и выходная проекция:
\[
z_t=W_{\mathrm{out}}h_t+b_{\mathrm{out}},
\qquad
z_t\in\R^V.
\]
Вектор $z_t$ называется вектором \emph{логитов}. Вероятности:
\[
p_\theta(x_{t+1}=v\mid x_{\leq t})
=
\frac{\exp(z_{t,v})}
{\sum_{u\in\V}\exp(z_{t,u})}.
\]

Матрицу $W_{\mathrm{out}}$ часто связывают с матрицей входных
эмбеддингов:
\[
W_{\mathrm{out}}=E^\top.
\]
Это называется \emph{weight tying}.

\subsection{Основные семейства Transformer}

\begin{center}
\begin{tabularx}{\textwidth}{>{\bfseries}l X X}
\toprule
Архитектура & Схема внимания & Типичные задачи \\
\midrule
Encoder-only
& Двунаправленное внимание ко всей входной последовательности
& Классификация, извлечение признаков, поиск, разметка токенов
\\
Decoder-only
& Каузальное внимание слева направо
& Продолжение текста, диалог, программирование, универсальная генерация
\\
Encoder--decoder
& Двунаправленный encoder и каузальный decoder с cross-attention
& Перевод, суммаризация, условная генерация
\\
\bottomrule
\end{tabularx}
\end{center}

В cross-attention запросы поступают от decoder, а ключи и значения ---
от encoder:
\[
Q=H_{\mathrm{dec}}W_Q,
\qquad
K=H_{\mathrm{enc}}W_K,
\qquad
V=H_{\mathrm{enc}}W_V.
\]

\subsection{Оценка числа параметров}

Для decoder-only модели со стандартным MHA и обычным двухслойным FFN:

\begin{itemize}
    \item эмбеддинги: приблизительно $Vd$;
    \item внимание одного слоя: приблизительно $4d^2$;
    \item FFN одного слоя: приблизительно $2dd_{\mathrm{ff}}$.
\end{itemize}

При связанном выходном слое:
\[
N_{\mathrm{params}}
\approx
Vd+
L\left(4d^2+2dd_{\mathrm{ff}}\right).
\]

\begin{example}
Пусть
\[
V=32000,\qquad
d=768,\qquad
d_{\mathrm{ff}}=3072,\qquad
L=12.
\]
Тогда
\[
Vd=24\,576\,000.
\]
Один слой:
\[
4d^2+2dd_{\mathrm{ff}}
=
4\cdot768^2+2\cdot768\cdot3072
=
7\,077\,888.
\]
Двенадцать слоёв:
\[
12\cdot7\,077\,888=84\,934\,656.
\]
Всего без малых параметров нормализаций и смещений:
\[
N_{\mathrm{params}}\approx109.5\text{ млн}.
\]
Если выходная матрица не связана с эмбеддингами, дополнительно требуется
ещё $Vd$ параметров.
\end{example}

\subsection{Вычислительная сложность}

Для последовательности длины $n$:

\begin{itemize}
    \item формирование $Q,K,V$ имеет порядок $O(nd^2)$;
    \item матрица $QK^\top$ --- $O(n^2d)$;
    \item умножение матрицы внимания на $V$ --- $O(n^2d)$;
    \item FFN --- $O(ndd_{\mathrm{ff}})$.
\end{itemize}

Квадратичная зависимость внимания от $n$ ограничивает работу с длинным
контекстом. На практике существенны не только арифметические операции,
но и чтение и запись данных в память ускорителя.

\subsection{Mixture of Experts}

В Mixture of Experts плотный FFN заменяется набором экспертов
$E_1,\ldots,E_m$. Маршрутизатор вычисляет оценки
\[
g(x)=W_rx
\]
и выбирает один или несколько экспертов:
\[
y=\sum_{j\in\operatorname{TopK}(g(x))}
\alpha_jE_j(x).
\]

Основные шаги:

\begin{enumerate}
    \item вычислить оценки маршрутизации токена;
    \item выбрать top-$k$ экспертов;
    \item передать токен выбранным экспертам;
    \item объединить их выходы;
    \item использовать вспомогательный штраф для равномерной загрузки.
\end{enumerate}

Общее число параметров может быть большим, но для одного токена
активируется только малая часть. Цена этого подхода --- сложная
маршрутизация, межустройственная коммуникация и возможная
несбалансированность экспертов.

% ============================================================
\section{Данные и проектирование модели}
% ============================================================

\subsection{Выбор конфигурации}

До обучения задаются:

\begin{itemize}
    \item размер словаря $V$;
    \item размерность $d$;
    \item число слоёв $L$;
    \item число голов запросов и голов ключей/значений;
    \item размер FFN;
    \item максимальная длина контекста;
    \item тип позиционного кодирования;
    \item тип нормализации и функции активации;
    \item плотная или разреженная архитектура;
    \item объём данных и вычислительный бюджет.
\end{itemize}

Увеличение числа параметров без достаточного количества данных ведёт к
недообученной модели. Увеличение данных при слишком малой модели может
упереться в ограниченную ёмкость сети.

\subsection{Эмпирические законы масштабирования}

Наблюдаемую зависимость потерь от числа параметров $N$ и числа обучающих
токенов $D$ часто аппроксимируют степенным выражением
\[
\mathcal{L}(N,D)
\approx
\mathcal{L}_{\infty}
+
\frac{A}{N^\alpha}
+
\frac{B}{D^\beta}.
\]

Это не математический закон, действующий для любых данных и архитектур,
а эмпирическая модель. Коэффициенты зависят от качества корпуса,
токенизации, архитектуры и режима обучения.

Грубая инженерная оценка вычислительной стоимости обучения плотного
Transformer:
\[
C\approx 6ND
\]
операций с плавающей точкой, где $N$ --- число параметров, а $D$ ---
число обработанных токенов. Коэффициент является приближённым и зависит
от того, какие операции учитываются.

Compute-optimal проектирование требует совместно выбирать $N$ и $D$:
вычислительный бюджет следует распределять не только на увеличение
модели, но и на достаточное число качественных токенов.

\subsection{Источники данных}

Корпус может содержать:

\begin{itemize}
    \item веб-страницы;
    \item книги и статьи;
    \item энциклопедические тексты;
    \item научные публикации;
    \item исходный код;
    \item документацию;
    \item диалоги;
    \item специализированные предметные данные;
    \item синтетические данные.
\end{itemize}

Для каждого источника необходимо учитывать:

\begin{itemize}
    \item происхождение и разрешённость использования;
    \item лицензионные условия;
    \item конфиденциальность и персональные данные;
    \item язык и предметную область;
    \item качество и достоверность;
    \item риск вредного или нежелательного содержания;
    \item риск пересечения с тестовыми наборами.
\end{itemize}

\subsection{Конвейер подготовки данных}

Типичный конвейер:

\begin{enumerate}
    \item загрузка и фиксация происхождения документа;
    \item извлечение основного текста;
    \item удаление шаблонной разметки, меню и навигации;
    \item определение языка;
    \item фильтрация повреждённых и бессодержательных документов;
    \item поиск точных дубликатов;
    \item поиск близких дубликатов;
    \item удаление тестовых и запрещённых данных;
    \item поиск персональной или секретной информации;
    \item оценка качества;
    \item формирование смеси источников;
    \item токенизация и упаковка последовательностей.
\end{enumerate}

\subsection{Дедупликация}

Точные дубликаты можно находить по криптографическому хешу
нормализованного документа.

Для близких дубликатов применяется схема на основе шинглов:

\begin{enumerate}
    \item разбить документ на символьные или токенные $k$-граммы;
    \item представить документ множеством шинглов;
    \item вычислить компактную MinHash-сигнатуру;
    \item с помощью LSH найти кандидатов с похожими сигнатурами;
    \item для кандидатов вычислить более точную меру сходства;
    \item оставить один репрезентативный документ из группы.
\end{enumerate}

Сходство Жаккара множеств шинглов $A$ и $B$:
\[
J(A,B)=\frac{|A\cap B|}{|A\cup B|}.
\]

Дедупликация уменьшает:

\begin{itemize}
    \item переобучение на повторяющихся фрагментах;
    \item запоминание текста;
    \item смещение частот;
    \item утечки между обучающей и тестовой выборками;
    \item расход вычислений на повторные примеры.
\end{itemize}

\subsection{Смешивание источников}

Если источник $i$ содержит $n_i$ токенов, его вероятность выборки можно
задать как
\[
p_i=
\frac{n_i^\alpha}
{\sum_j n_j^\alpha},
\qquad 0<\alpha\leq 1.
\]

При $\alpha=1$ источники выбираются пропорционально размеру. При
$\alpha<1$ небольшие источники относительно усиливаются. Это полезно для
качественных, специализированных и малоресурсных данных, но слишком
сильное повышение веса приводит к многократному повторению одних и тех
же документов.

\subsection{Упаковка последовательностей}

Документы объединяют в блоки фиксированной длины $T$:
\[
[x^{(1)},\texttt{<EOS>},x^{(2)},\texttt{<EOS>},\ldots].
\]

Упаковка уменьшает количество $\texttt{<PAD>}$ и повышает использование
ускорителя. При этом нужно решить, может ли внимание переходить через
границу документов. Возможны два режима:

\begin{enumerate}
    \item единая каузальная маска, позволяющая видеть предыдущий документ;
    \item блочно-диагональная маска, изолирующая документы.
\end{enumerate}

\subsection{Контроль загрязнения тестов}

Если тестовые задания попали в обучающий корпус, результат может
измерять запоминание, а не обобщение.

Практические меры:

\begin{itemize}
    \item зафиксировать тесты до окончательной подготовки данных;
    \item удалять точные совпадения;
    \item искать совпадения по длинным $n$-граммам;
    \item искать переформулированные и близкие варианты;
    \item проверять даты появления тестов и документов;
    \item отдельно сообщать результаты на закрытых и новых тестах.
\end{itemize}

% ============================================================
\section{Предварительное обучение}
% ============================================================

\subsection{Основной алгоритм}

\begin{algorithm}[h]
\caption{Предварительное обучение авторегрессионной LLM}
\begin{algorithmic}[1]
\Require токенизированный корпус $\D$; модель $p_\theta$;
число шагов $S$
\Ensure обученные параметры $\theta$
\For{$s=1,\ldots,S$}
    \State выбрать пакет токенных блоков $B\in\mathbb{N}^{b\times T}$
    \State сформировать входы $X=B[:,1:T-1]$
    \State сформировать цели $Y=B[:,2:T]$
    \State вычислить логиты $Z=\Call{Transformer}{X;\theta}$
    \State вычислить среднюю кросс-энтропию
    \[
    \mathcal{L}
    =
    -\frac{1}{N}
    \sum_{i,t}
    \log p_\theta(Y_{i,t}\mid X_{i,\leq t})
    \]
    \State выполнить обратное распространение ошибки
    \State при необходимости обрезать норму градиента
    \State обновить параметры оптимизатором
    \State обновить расписание скорости обучения
    \State сохранить метрики и периодический checkpoint
\EndFor
\State \Return $\theta$
\end{algorithmic}
\end{algorithm}

\subsection{Оптимизатор AdamW}

Пусть на шаге $t$
\[
g_t=\nabla_\theta\mathcal{L}_t(\theta_{t-1}).
\]
Оценки первого и второго моментов:
\[
m_t=\beta_1m_{t-1}+(1-\beta_1)g_t,
\]
\[
v_t=\beta_2v_{t-1}+(1-\beta_2)g_t^2.
\]
Коррекция смещения:
\[
\widehat{m}_t=\frac{m_t}{1-\beta_1^t},
\qquad
\widehat{v}_t=\frac{v_t}{1-\beta_2^t}.
\]
Обновление AdamW:
\[
\theta_t
=
(1-\eta_t\lambda)\theta_{t-1}
-
\eta_t
\frac{\widehat{m}_t}
{\sqrt{\widehat{v}_t}+\varepsilon}.
\]

Здесь:

\begin{itemize}
    \item $\eta_t$ --- скорость обучения;
    \item $\lambda$ --- коэффициент weight decay;
    \item $\beta_1,\beta_2$ --- коэффициенты экспоненциального сглаживания.
\end{itemize}

В AdamW weight decay отделён от градиента функции потерь. Это отличает
его от добавления обычного $L_2$-штрафа внутрь адаптивного оптимизатора.

\subsection{Расписание скорости обучения}

Часто применяется:

\begin{enumerate}
    \item \textbf{warm-up}: линейное увеличение $\eta_t$ от малого
    значения;
    \item \textbf{основная стадия}: постоянная или медленно меняющаяся
    скорость;
    \item \textbf{decay}: косинусное или линейное уменьшение.
\end{enumerate}

Косинусное расписание после warm-up:
\[
\eta_t
=
\eta_{\min}
+
\frac{\eta_{\max}-\eta_{\min}}{2}
\left(
1+\cos\frac{\pi(t-t_{\mathrm{warm}})}
{T-t_{\mathrm{warm}}}
\right).
\]

Warm-up уменьшает риск нестабильности в начале, когда оценки моментов
оптимизатора и масштабы активаций ещё не установились.

\subsection{Стабилизация}

Основные техники:

\begin{itemize}
    \item Pre-Norm блоки;
    \item нормализация входных данных и активаций;
    \item gradient clipping:
    \[
    g\leftarrow
    g\cdot
    \min\left(1,\frac{c}{\|g\|_2}\right);
    \]
    \item warm-up;
    \item корректная инициализация;
    \item контроль выбросов в данных;
    \item BF16 или смешанная точность;
    \item накопление градиента;
    \item периодические checkpoints;
    \item автоматический перезапуск после сбоев.
\end{itemize}

\subsection{Смешанная точность}

Часть вычислений выполняется в FP16, BF16 или FP8, а критичные состояния
могут храниться с большей точностью.

Преимущества:

\begin{itemize}
    \item меньший расход памяти;
    \item большая пропускная способность;
    \item использование специализированных матричных блоков ускорителя.
\end{itemize}

Риски:

\begin{itemize}
    \item переполнение и потеря малых значений;
    \item нестабильность softmax и нормализации;
    \item необходимость loss scaling для FP16;
    \item зависимость от реализации отдельных операций.
\end{itemize}

BF16 имеет более широкий диапазон экспоненты, чем FP16, и поэтому часто
удобнее для обучения больших моделей.

\subsection{Параллельное обучение}

\begin{description}
    \item[Data parallelism.] Каждое устройство содержит копию модели,
    но получает свою часть пакета. Градиенты объединяются операцией
    all-reduce.

    \item[Tensor parallelism.] Отдельные матрицы и операции одного слоя
    разделяются между устройствами.

    \item[Pipeline parallelism.] Последовательные группы слоёв размещаются
    на разных устройствах; пакет делится на микропакеты.

    \item[Sequence/context parallelism.] Токенные позиции или часть
    промежуточных состояний разделяются между устройствами.

    \item[Expert parallelism.] Эксперты MoE распределяются между
    устройствами.
\end{description}

При data parallelism полный набор параметров, градиентов и состояний
оптимизатора может не помещаться в память. ZeRO-подобные методы
разделяют между устройствами:

\begin{enumerate}
    \item состояния оптимизатора;
    \item градиенты;
    \item параметры.
\end{enumerate}

\subsection{Activation checkpointing}

Вместо хранения всех активаций сохраняется только часть. Недостающие
активации повторно вычисляются при обратном проходе.

Это уменьшает память ценой дополнительных вычислений:
\[
\text{меньше памяти}
\quad\Longleftrightarrow\quad
\text{больше повторных вычислений}.
\]

\subsection{Эффективный размер пакета}

Если используются:

\begin{itemize}
    \item микропакет размера $b$;
    \item $a$ шагов накопления градиента;
    \item $g$ data-parallel устройств,
\end{itemize}

то эффективный размер пакета:
\[
B_{\mathrm{eff}}=b\cdot a\cdot g.
\]

Для языковых моделей полезно также считать число токенов в пакете:
\[
T_{\mathrm{batch}}
=
B_{\mathrm{eff}}\cdot T.
\]

\subsection{Мониторинг обучения}

Контролируются:

\begin{itemize}
    \item обучающая и валидационная потеря;
    \item норма градиента;
    \item скорость обучения;
    \item доля пропущенных или повреждённых пакетов;
    \item скорость обработки токенов;
    \item загрузка устройств;
    \item объём коммуникаций;
    \item частота переполнений;
    \item качество на промежуточных тестах;
    \item распределение данных по источникам.
\end{itemize}

Резкий всплеск loss может быть вызван повреждённым документом,
слишком высокой скоростью обучения, численной нестабильностью или
ошибкой распределённой синхронизации.

% ============================================================
\section{Последующее обучение и адаптация}
% ============================================================

\subsection{Continued Pretraining}

Продолженное предварительное обучение использует ту же
авторегрессионную функцию потерь, но на новом корпусе.

Варианты:

\begin{description}
    \item[Domain-adaptive pretraining.] Адаптация к предметной области.
    \item[Language-adaptive pretraining.] Усиление конкретного языка.
    \item[Long-context adaptation.] Дополнительное обучение на длинных
    последовательностях.
    \item[Continual pretraining.] Обновление знаний новыми документами.
\end{description}

Чтобы уменьшить катастрофическое забывание, специализированные данные
смешивают с частью исходных общих данных.

\subsection{Supervised Fine-Tuning}

Пусть имеется набор
\[
\D_{\mathrm{SFT}}=\{(x^{(i)},y^{(i)})\},
\]
где $x$ --- инструкция и контекст, а $y$ --- эталонный ответ.

Функция потерь:
\[
\mathcal{L}_{\mathrm{SFT}}
=
-\sum_i\sum_{t=1}^{|y^{(i)}|}
\log
\pi_\theta
\left(
y_t^{(i)}
\mid
x^{(i)},y_{<t}^{(i)}
\right).
\]

Обычно loss вычисляется только по токенам ответа ассистента:
\[
\mathcal{L}
=
-\frac{1}{\sum_t m_t}
\sum_t
m_t\log p_\theta(x_t\mid x_{<t}),
\]
где
\[
m_t=
\begin{cases}
1,&\text{если токен принадлежит целевому ответу},\\
0,&\text{если токен принадлежит инструкции или служебной части}.
\end{cases}
\]

Качество SFT-данных часто важнее их абсолютного количества. Данные должны
содержать корректные инструкции, разнообразные задачи, отказ от
невыполнимых запросов и требуемый формат ответа.

\subsection{Обучение по предпочтениям}

Пусть для запроса $x$ даны два ответа:

\[
y_w \succ y_l,
\]
где $y_w$ предпочтительнее $y_l$.

Reward model обучается присваивать ответу скалярную оценку $r_\phi(x,y)$.
Модель Брэдли--Терри задаёт
\[
\Pr(y_w\succ y_l\mid x)
=
\sigma
\left(
r_\phi(x,y_w)-r_\phi(x,y_l)
\right),
\]
где
\[
\sigma(z)=\frac{1}{1+e^{-z}}.
\]

Функция потерь reward model:
\[
\mathcal{L}_{\mathrm{RM}}
=
-\E_{(x,y_w,y_l)}
\log\sigma
\left(
r_\phi(x,y_w)-r_\phi(x,y_l)
\right).
\]

\subsection{RLHF}

Классический конвейер Reinforcement Learning from Human Feedback:

\begin{enumerate}
    \item обучить базовую модель;
    \item выполнить SFT на демонстрациях;
    \item сгенерировать несколько ответов на каждый запрос;
    \item собрать человеческие ранжирования;
    \item обучить reward model;
    \item оптимизировать языковую модель по предсказанной награде;
    \item ограничивать отклонение от опорной модели KL-штрафом;
    \item повторно оценивать полезность и безопасность.
\end{enumerate}

Типичная целевая функция:
\[
\max_\pi
\E_{y\sim\pi(\cdot\mid x)}
[r_\phi(x,y)]
-
\beta
\KL
\left(
\pi(\cdot\mid x)
\Vert
\pi_{\mathrm{ref}}(\cdot\mid x)
\right).
\]

KL-штраф предотвращает слишком сильное отклонение от исходной
инструктивной модели и уменьшает эксплуатацию ошибок reward model.

\subsection{Оптимальная политика при KL-регуляризации}

\begin{theorem}
Пусть для фиксированного запроса $x$ множество возможных ответов
$\mathcal{Y}$ конечно, $\pi_{\mathrm{ref}}(y)>0$ для всех
$y\in\mathcal{Y}$, $r(y)\in\R$, $\beta>0$. Рассмотрим задачу
\[
\max_{\pi}
\left[
\sum_{y\in\mathcal{Y}}\pi(y)r(y)
-
\beta
\sum_{y\in\mathcal{Y}}
\pi(y)\log\frac{\pi(y)}{\pi_{\mathrm{ref}}(y)}
\right]
\]
при ограничениях
\[
\pi(y)\geq0,\qquad
\sum_y\pi(y)=1.
\]

Единственное оптимальное распределение равно
\[
\pi^\ast(y)
=
\frac{
\pi_{\mathrm{ref}}(y)\exp(r(y)/\beta)
}{
Z
},
\]
где
\[
Z=
\sum_{z\in\mathcal{Y}}
\pi_{\mathrm{ref}}(z)\exp(r(z)/\beta).
\]
\end{theorem}

\begin{proof}
Введём множитель Лагранжа $\lambda$ для ограничения нормировки:
\[
\mathcal{F}(\pi,\lambda)
=
\sum_y\pi(y)r(y)
-
\beta\sum_y
\pi(y)\log\frac{\pi(y)}{\pi_{\mathrm{ref}}(y)}
+
\lambda\left(\sum_y\pi(y)-1\right).
\]

Для внутренней точки симплекса:
\begin{align*}
\frac{\partial\mathcal{F}}{\partial\pi(y)}
&=
r(y)
-
\beta
\left(
\log\frac{\pi(y)}{\pi_{\mathrm{ref}}(y)}+1
\right)
+
\lambda.
\end{align*}
В стационарной точке производная равна нулю:
\[
r(y)
-
\beta
\left(
\log\frac{\pi(y)}{\pi_{\mathrm{ref}}(y)}+1
\right)
+
\lambda=0.
\]
Следовательно,
\[
\log\frac{\pi(y)}{\pi_{\mathrm{ref}}(y)}
=
\frac{r(y)}{\beta}
+
\frac{\lambda}{\beta}
-1.
\]
После экспоненцирования:
\[
\pi(y)
=
\pi_{\mathrm{ref}}(y)
\exp(r(y)/\beta)
\exp(\lambda/\beta-1).
\]
Последний множитель не зависит от $y$. Из условия
$\sum_y\pi(y)=1$ получаем
\[
\exp(\lambda/\beta-1)=\frac{1}{Z}.
\]
Поэтому
\[
\pi^\ast(y)
=
\frac{
\pi_{\mathrm{ref}}(y)\exp(r(y)/\beta)
}{
Z
}.
\]

Линейный член по $\pi$ является одновременно вогнутым и выпуклым, а
$-\KL(\pi\Vert\pi_{\mathrm{ref}})$ строго вогнут по $\pi$ на внутренности
симплекса. Следовательно, вся целевая функция строго вогнута и найденная
стационарная точка является единственным глобальным максимумом.
\end{proof}

Из формулы следует
\[
r(y)
=
\beta
\log\frac{\pi^\ast(y)}{\pi_{\mathrm{ref}}(y)}
+
\beta\log Z.
\]
Для разности наград нормировочная константа сокращается:
\[
r(y_w)-r(y_l)
=
\beta
\left[
\log\frac{\pi^\ast(y_w)}{\pi_{\mathrm{ref}}(y_w)}
-
\log\frac{\pi^\ast(y_l)}{\pi_{\mathrm{ref}}(y_l)}
\right].
\]

\subsection{Direct Preference Optimization}

DPO непосредственно оптимизирует языковую модель по парам предпочтений,
не обучая отдельную reward model и не запуская PPO.

Функция потерь:
\[
\mathcal{L}_{\mathrm{DPO}}
=
-\E_{(x,y_w,y_l)}
\log\sigma
\left(
\beta
\left[
\log
\frac{\pi_\theta(y_w\mid x)}
{\pi_{\mathrm{ref}}(y_w\mid x)}
-
\log
\frac{\pi_\theta(y_l\mid x)}
{\pi_{\mathrm{ref}}(y_l\mid x)}
\right]
\right).
\]

Алгоритм:

\begin{enumerate}
    \item зафиксировать опорную модель $\pi_{\mathrm{ref}}$;
    \item для каждой пары вычислить логарифмы вероятностей выбранного и
    отклонённого ответов;
    \item вычесть соответствующие логарифмы вероятностей опорной модели;
    \item вычислить логистическую loss;
    \item обновить только обучаемую модель $\pi_\theta$.
\end{enumerate}

DPO проще классического RLHF, но по-прежнему зависит от качества и
репрезентативности предпочтений.

\subsection{LoRA}

Пусть в исходной модели имеется линейное преобразование
\[
y=Wx,
\qquad
W\in\R^{d_{\mathrm{out}}\times d_{\mathrm{in}}}.
\]
В LoRA матрица $W$ замораживается, а обновление представляется
низкоранговым произведением:
\[
W'=W+\frac{\alpha}{r}BA,
\]
где
\[
A\in\R^{r\times d_{\mathrm{in}}},
\qquad
B\in\R^{d_{\mathrm{out}}\times r},
\qquad
r\ll
\min(d_{\mathrm{in}},d_{\mathrm{out}}).
\]

Обучаются только $A$ и $B$. Число обучаемых параметров:
\[
r(d_{\mathrm{in}}+d_{\mathrm{out}})
\]
вместо
\[
d_{\mathrm{in}}d_{\mathrm{out}}.
\]

\begin{example}
Для
\[
W\in\R^{4096\times4096},
\qquad r=8
\]
полное обучение требует
\[
4096^2=16\,777\,216
\]
обучаемых параметров.

LoRA требует
\[
8(4096+4096)=65\,536
\]
параметров, то есть в $256$ раз меньше для этой матрицы.
\end{example}

После обучения LoRA-обновление можно:

\begin{itemize}
    \item хранить отдельно как адаптер;
    \item динамически подключать к базовой модели;
    \item объединить с $W$ перед инференсом.
\end{itemize}

\subsection{QLoRA}

QLoRA сочетает:

\begin{enumerate}
    \item замороженную квантованную базовую модель;
    \item вычисление градиента через операции деквантования;
    \item обучаемые LoRA-адаптеры с более высокой точностью.
\end{enumerate}

Это значительно уменьшает память для fine-tuning. Квантованные базовые
веса не обновляются, а новая информация записывается в адаптеры.

\subsection{Выбор способа адаптации}

\begin{center}
\begin{tabularx}{\textwidth}{>{\bfseries}l X}
\toprule
Задача & Предпочтительный подход \\
\midrule
Изменить формат, стиль и выполнение инструкций
& SFT или LoRA
\\
Учесть человеческие предпочтения
& DPO или RLHF
\\
Добавить обновляемые факты
& RAG
\\
Адаптировать язык и терминологию
& Continued pretraining, затем SFT
\\
Сохранить несколько дешёвых специализаций
& Отдельные LoRA-адаптеры
\\
Нужна максимальная предметная адаптация
& Полный fine-tuning при наличии ресурсов и данных
\\
\bottomrule
\end{tabularx}
\end{center}

% ============================================================
\section{Авторегрессионный инференс}
% ============================================================

\subsection{Базовый алгоритм генерации}

\begin{algorithm}[h]
\caption{Авторегрессионная генерация}
\begin{algorithmic}[1]
\Require входные токены $x_{1:m}$; максимальное число новых токенов $K$
\Ensure продолженная последовательность
\State $s\gets[x_1,\ldots,x_m]$
\For{$k=1,\ldots,K$}
    \State вычислить логиты $z=\Call{Model}{s}$
    \State взять логиты последней позиции $z_{\mathrm{last}}$
    \State применить температурное масштабирование и ограничения
    \State получить распределение $p$ с помощью softmax
    \State выбрать или случайно сэмплировать следующий токен $v$
    \State добавить $v$ в конец $s$
    \If{$v=\texttt{<EOS>}$ или сработало условие остановки}
        \State \textbf{break}
    \EndIf
\EndFor
\State \Return $s$
\end{algorithmic}
\end{algorithm}

Без KV-кэша на каждом шаге пришлось бы повторно вычислять ключи и значения
всего префикса.

\subsection{Температура}

Для $T>0$:
\[
p_i(T)
=
\frac{\exp(z_i/T)}
{\sum_j\exp(z_j/T)}.
\]

\begin{itemize}
    \item $T<1$ делает распределение более концентрированным;
    \item $T>1$ делает его более плоским;
    \item предел $T\to0$ соответствует выбору максимального логита;
    \item температура не добавляет модели знаний, а только изменяет
    случайность выбора.
\end{itemize}

\subsection{Greedy decoding}

Выбирается наиболее вероятный токен:
\[
x_{t+1}
=
\argmax_{v\in\V}
p_\theta(v\mid x_{\leq t}).
\]

Преимущества:

\begin{itemize}
    \item детерминированность;
    \item низкая вычислительная стоимость;
    \item удобство для строгих задач.
\end{itemize}

Недостатки:

\begin{itemize}
    \item локально лучший выбор не обязан давать лучшую последовательность;
    \item возможны повторения;
    \item низкое разнообразие.
\end{itemize}

\subsection{Top-$k$ sampling}

Сохраняются $k$ токенов с наибольшими вероятностями. Остальные получают
нулевую вероятность, после чего распределение нормируется.

Пусть
\[
p=(0.50,0.30,0.15,0.05).
\]
При $k=2$:
\[
\widetilde{p}
=
\left(
\frac{0.50}{0.80},
\frac{0.30}{0.80},
0,
0
\right)
=
(0.625,0.375,0,0).
\]

\subsection{Nucleus, или top-$p$, sampling}

Токены сортируются по убыванию вероятности. Выбирается минимальное
множество $S$, для которого
\[
\sum_{v\in S}p(v)\geq p_0.
\]
После этого вероятности вне $S$ обнуляются и выполняется нормировка.

В предыдущем примере при $p_0=0.8$ остаются первые два токена.

В отличие от top-$k$, число допустимых токенов адаптируется к форме
распределения.

\subsection{Beam search}

Beam search хранит $B$ лучших частичных последовательностей.

Основные шаги:

\begin{enumerate}
    \item начать с одной пустой гипотезы;
    \item расширить каждую гипотезу допустимыми токенами;
    \item вычислить сумму логарифмов вероятностей;
    \item оставить $B$ лучших продолжений;
    \item повторять до завершения;
    \item выбрать лучшую законченную последовательность.
\end{enumerate}

Для уменьшения предпочтения коротких ответов применяют нормировку:
\[
\operatorname{score}(y)
=
\frac{\log p(y\mid x)}
{|y|^\alpha}.
\]

Beam search полезен для перевода и задач с узким пространством ответов,
но для открытого диалога может давать однообразный текст.

\subsection{Управление повторениями}

Применяются:

\begin{itemize}
    \item штраф за повтор токена;
    \item frequency penalty;
    \item presence penalty;
    \item запрет повторяющихся $n$-грамм;
    \item специальные stop-последовательности;
    \item ограничение максимальной длины.
\end{itemize}

Эти методы являются эвристиками. Слишком сильные штрафы могут ломать
код, формулы, имена и естественные повторы.

\subsection{KV-кэш}

На каждом слое сохраняются ключи и значения предыдущих токенов:
\[
K_{1:t}^{(\ell)},\qquad
V_{1:t}^{(\ell)}.
\]
При генерации нового токена вычисляются только
\[
q_{t+1}^{(\ell)},\qquad
k_{t+1}^{(\ell)},\qquad
v_{t+1}^{(\ell)},
\]
а новые ключ и значение добавляются к кэшу.

Приближённый размер KV-кэша на одну последовательность:
\[
M_{\mathrm{KV}}
\approx
2Lnh_{\mathrm{kv}}d_{\mathrm{head}}b,
\]
где:

\begin{itemize}
    \item множитель $2$ соответствует ключам и значениям;
    \item $L$ --- число слоёв;
    \item $n$ --- длина контекста;
    \item $h_{\mathrm{kv}}$ --- число голов ключей/значений;
    \item $d_{\mathrm{head}}$ --- размер головы;
    \item $b$ --- число байтов на элемент.
\end{itemize}

\begin{example}
Пусть
\[
L=32,\quad n=4096,\quad
h_{\mathrm{kv}}=8,\quad
d_{\mathrm{head}}=128,\quad b=2.
\]
Тогда
\[
M_{\mathrm{KV}}
=
2\cdot32\cdot4096\cdot8\cdot128\cdot2
=
536\,870\,912
\]
байт, то есть примерно $512$ MiB на одну последовательность.
\end{example}

GQA и MQA уменьшают $h_{\mathrm{kv}}$ и тем самым сокращают KV-кэш.

\subsection{Prefill и decode}

Инференс делится на две стадии:

\begin{description}
    \item[Prefill.] Весь исходный запрос обрабатывается параллельно,
    формируется KV-кэш.

    \item[Decode.] Новые токены генерируются последовательно по одному.
\end{description}

Prefill обычно лучше использует вычислительные блоки, а decode часто
ограничен пропускной способностью памяти, поскольку веса модели нужно
многократно читать для генерации небольшого числа новых активаций.

\subsection{Оптимизация инференса}

\begin{description}
    \item[Continuous batching.] Запросы добавляются в динамический пакет
    по мере освобождения мест.

    \item[Paged KV-cache.] KV-кэш хранится блоками, что уменьшает
    фрагментацию памяти.

    \item[Prefix caching.] Общий префикс нескольких запросов вычисляется
    один раз.

    \item[FlashAttention.] Вычисляет точное внимание блоками, не сохраняя
    полностью промежуточную матрицу размера $n\times n$ во внешней памяти.

    \item[Квантование.] Веса и иногда активации хранятся с меньшим числом
    бит.

    \item[Speculative decoding.] Малая модель предлагает несколько
    токенов, большая проверяет их параллельно.

    \item[Разделение prefill и decode.] Стадии могут обслуживаться
    разными группами ускорителей.
\end{description}

\subsection{Квантование}

Равномерное квантование приближает число $w$:
\[
q=\operatorname{round}\left(\frac{w}{s}\right)+z,
\]
где $s$ --- масштаб, $z$ --- нулевая точка. Деквантованное значение:
\[
\widehat{w}=s(q-z).
\]

Возможны:

\begin{itemize}
    \item weight-only quantization;
    \item квантование весов и активаций;
    \item квантование KV-кэша;
    \item поканальное и погрупповое квантование;
    \item статическая или динамическая калибровка.
\end{itemize}

Меньшая точность уменьшает память и трафик, но может ухудшать качество,
особенно для выбросов и чувствительных слоёв.

\subsection{Speculative decoding}

Пусть $q$ --- распределение малой draft-модели, а $p$ ---
распределение целевой модели.

Один шаг точной проверки:

\begin{enumerate}
    \item draft-модель предлагает токен $y\sim q$;
    \item токен принимается с вероятностью
    \[
    a(y)=\min\left(1,\frac{p(y)}{q(y)}\right);
    \]
    \item при отклонении новый токен выбирается из
    \[
    r(v)
    =
    \frac{(p(v)-q(v))_+}
    {\sum_u(p(u)-q(u))_+}.
    \]
\end{enumerate}

Проверим итоговое распределение. Вероятность получить $y$ через принятие:
\[
q(y)a(y)=\min(p(y),q(y)).
\]
Вероятность отклонения:
\[
R
=
1-\sum_y\min(p(y),q(y))
=
\sum_y(p(y)-q(y))_+.
\]
Вклад ветви отклонения:
\[
Rr(y)=(p(y)-q(y))_+.
\]
Итог:
\[
\min(p(y),q(y))+(p(y)-q(y))_+=p(y).
\]
Следовательно, корректирующая процедура сохраняет распределение целевой
модели.

Для нескольких предложенных токенов большая модель проверяет блок за
один параллельный forward pass. Ускорение зависит от качества
draft-модели и доли принятых токенов.

% ============================================================
\section{Практическое использование LLM}
% ============================================================

\subsection{Выбор модели}

Следует учитывать:

\begin{itemize}
    \item качество на целевой задаче;
    \item поддерживаемые языки;
    \item размер контекста;
    \item фактическую способность использовать длинный контекст;
    \item задержку и пропускную способность;
    \item стоимость;
    \item лицензию;
    \item возможность локального развёртывания;
    \item требования к конфиденциальности;
    \item поддержку инструментов и структурированного вывода;
    \item доступность fine-tuning;
    \item качество документации и воспроизводимость.
\end{itemize}

Наиболее крупная модель не всегда является оптимальной. Для простой,
частой и строго формализованной операции меньшая специализированная
модель может быть быстрее, дешевле и надёжнее.

\subsection{Структура промпта}

Хороший запрос обычно содержит:

\begin{enumerate}
    \item \textbf{задачу} --- что именно требуется;
    \item \textbf{контекст} --- данные, на которых нужно работать;
    \item \textbf{ограничения} --- язык, длина, запреты, критерии;
    \item \textbf{формат вывода} --- таблица, JSON, список, формула;
    \item \textbf{примеры} --- если формат трудно описать словами;
    \item \textbf{правило неопределённости} --- что делать при недостатке
    данных;
    \item \textbf{критерии проверки} --- тесты, ссылки, вычисления.
\end{enumerate}

Минимальный шаблон:

\begin{quote}
\textbf{Задача:} извлечь из текста название метода, вход и выход.

\textbf{Контекст:} текст между явными разделителями.

\textbf{Ограничения:} не использовать внешние факты; если поле отсутствует,
вернуть \texttt{null}.

\textbf{Формат:} JSON с полями
\texttt{method}, \texttt{input}, \texttt{output}.
\end{quote}

\subsection{Zero-shot, one-shot и few-shot}

\begin{description}
    \item[Zero-shot.] Передаётся только инструкция.
    \item[One-shot.] Добавляется один пример.
    \item[Few-shot.] Добавляется несколько примеров входа и ответа.
\end{description}

Примеры особенно полезны, когда важны:

\begin{itemize}
    \item нестандартная разметка;
    \item классификационные метки;
    \item стиль;
    \item граничные случаи;
    \item точная структура результата.
\end{itemize}

Примеры должны быть корректными и разнообразными. Модель может
воспроизводить ошибки, смещения и случайные особенности демонстраций.

\subsection{Декомпозиция задачи}

Сложную задачу полезно разделить:

\begin{enumerate}
    \item определить требуемые подзадачи;
    \item получить или вычислить промежуточные данные;
    \item проверить промежуточные результаты;
    \item сформировать итоговый ответ;
    \item выполнить независимую валидацию.
\end{enumerate}

В прикладных системах предпочтительны проверяемые промежуточные
артефакты: таблицы, формулы, тесты, вызовы инструментов и ссылки на
источники. Подробное текстовое рассуждение модели само по себе не
является гарантией правильности.

\subsection{Структурированный вывод}

Для JSON, SQL или вызова функции задаётся схема:

\begin{itemize}
    \item список полей;
    \item тип каждого поля;
    \item обязательность;
    \item допустимые значения;
    \item ограничения длины;
    \item правило обработки отсутствующих данных.
\end{itemize}

Результат необходимо валидировать программно. При ошибке можно:

\begin{enumerate}
    \item повторить декодирование с ограниченной грамматикой;
    \item вернуть модели сообщение об ошибке схемы;
    \item применить безопасное исправление;
    \item отклонить результат.
\end{enumerate}

\subsection{Retrieval-Augmented Generation}

RAG дополняет запрос фрагментами внешних документов.

Пусть документы разбиты на фрагменты
\[
c_1,\ldots,c_m.
\]
Для каждого фрагмента вычисляется вектор
\[
e_i=E_d(c_i).
\]
Для запроса:
\[
q=E_q(x).
\]
Поиск выполняется по косинусному сходству:
\[
s(q,e_i)
=
\frac{q^\top e_i}
{\|q\|_2\|e_i\|_2}.
\]

\begin{algorithm}[h]
\caption{Базовый RAG-конвейер}
\begin{algorithmic}[1]
\Require запрос $x$; индекс фрагментов $\{(c_i,e_i)\}$
\Ensure ответ $y$ и сведения об использованных фрагментах
\State вычислить вектор запроса $q=E_q(x)$
\State найти top-$k$ фрагментов по сходству
\State при необходимости выполнить reranking
\State отфильтровать документы по правам доступа и метаданным
\State собрать контекст $C$ из наиболее релевантных фрагментов
\State сформировать инструкцию отвечать только на основе $C$
\State сгенерировать ответ $y$
\State проверить, что ключевые утверждения поддержаны фрагментами
\State \Return $(y,C)$
\end{algorithmic}
\end{algorithm}

\begin{figure}[h]
\centering
\resizebox{0.95\textwidth}{!}{
\begin{tikzpicture}[node distance=12mm]
    \node[process] (query) {Запрос\\пользователя};
    \node[process, right=of query] (embed) {Векторизация\\запроса};
    \node[data, right=of embed] (index) {Векторный\\индекс};
    \node[process, right=of index] (retrieve) {Поиск и\\reranking};
    \node[process, right=of retrieve] (prompt) {Запрос +\\фрагменты};
    \node[process, right=of prompt] (llm) {LLM};
    \node[process, right=of llm] (answer) {Ответ и\\ссылки};

    \draw[arrow] (query) -- (embed);
    \draw[arrow] (embed) -- (index);
    \draw[arrow] (index) -- (retrieve);
    \draw[arrow] (retrieve) -- (prompt);
    \draw[arrow] (prompt) -- (llm);
    \draw[arrow] (llm) -- (answer);
\end{tikzpicture}
}
\caption{Схема Retrieval-Augmented Generation.}
\end{figure}

\begin{example}
В базе имеются фрагменты:

\begin{itemize}
    \item $c_1$: <<Поле \texttt{status} принимает значения
    \texttt{draft} и \texttt{approved}.>>
    \item $c_2$: <<Поле \texttt{priority} принимает значения от 1 до 5.>>
\end{itemize}

Запрос: <<Какие значения допустимы для \texttt{status}?>>

Retriever должен вернуть $c_1$, после чего модель формирует ответ
\texttt{draft} и \texttt{approved}. Если нужный фрагмент не найден,
система должна сообщить о недостатке данных, а не выдумывать значения.
\end{example}

\subsection{Проектирование RAG}

Качество зависит от всех стадий:

\begin{itemize}
    \item очистки документов;
    \item размера и перекрытия фрагментов;
    \item сохранения заголовков и метаданных;
    \item embedding-модели;
    \item векторного, лексического или гибридного поиска;
    \item reranker;
    \item числа фрагментов;
    \item порядка фрагментов;
    \item ограничения контекстного окна;
    \item инструкции по цитированию;
    \item проверки поддержки ответа источниками.
\end{itemize}

RAG не гарантирует истинность автоматически. Ошибка может возникнуть,
если:

\begin{enumerate}
    \item нужного документа нет в базе;
    \item документ плохо разбит;
    \item retriever выбрал нерелевантный фрагмент;
    \item релевантный фрагмент потерялся среди лишнего контекста;
    \item модель неверно интерпретировала источник;
    \item источник сам содержит ошибку.
\end{enumerate}

\subsection{Использование инструментов}

LLM может формировать структурированные вызовы:

\[
\text{название инструмента}
+
\text{аргументы}
\longrightarrow
\text{результат}.
\]

Примеры инструментов:

\begin{itemize}
    \item калькулятор;
    \item поиск;
    \item база данных;
    \item интерпретатор кода;
    \item календарь;
    \item система управления задачами;
    \item корпоративное API.
\end{itemize}

Безопасный цикл:

\begin{enumerate}
    \item модель предлагает действие;
    \item система проверяет имя инструмента;
    \item аргументы валидируются по схеме;
    \item проверяются права пользователя;
    \item опасное действие требует подтверждения;
    \item инструмент выполняется в изолированной среде;
    \item результат возвращается модели как данные;
    \item все действия журналируются;
    \item число шагов и время ограничиваются.
\end{enumerate}

\subsection{Агентный цикл}

Упрощённая схема агента:

\begin{enumerate}
    \item получить цель и текущее состояние;
    \item решить, достаточно ли данных для ответа;
    \item при необходимости выбрать инструмент;
    \item выполнить действие;
    \item получить наблюдение;
    \item обновить состояние;
    \item повторить до достижения цели или лимита;
    \item сформировать ответ и отчёт о выполненных действиях.
\end{enumerate}

Агент не должен самостоятельно получать неограниченные права. Следует
применять принцип наименьших привилегий.

\subsection{Защита от prompt injection}

Документы, веб-страницы и результаты инструментов являются
\emph{недоверенными данными}. Они могут содержать текст вида
<<игнорируй предыдущие инструкции>>.

Защитные меры:

\begin{itemize}
    \item разделять системные инструкции и внешние данные;
    \item явно сообщать, что текст документов не является инструкцией;
    \item не передавать модели секреты без необходимости;
    \item проверять права доступа до извлечения документа;
    \item ограничивать список инструментов;
    \item валидировать каждый вызов;
    \item требовать подтверждения необратимых операций;
    \item фильтровать вывод перед исполнением;
    \item использовать изоляцию и журналирование;
    \item не полагаться только на текстовый системный промпт.
\end{itemize}

% ============================================================
\section{Оценивание}
% ============================================================

\subsection{Уровни оценки}

\begin{enumerate}
    \item \textbf{Компонентный уровень}: tokenizer, retriever, reranker,
    tool selection.
    \item \textbf{Модельный уровень}: loss, accuracy, factuality,
    instruction following.
    \item \textbf{Системный уровень}: полный RAG или агентный конвейер.
    \item \textbf{Продуктовый уровень}: полезность, задержка, стоимость,
    безопасность.
\end{enumerate}

\subsection{Метрики качества}

\begin{description}
    \item[Accuracy.] Доля правильных ответов.
    \item[Exact Match.] Полное совпадение после нормализации.
    \item[Precision, Recall, $F_1$.] Используются для извлечения,
    классификации и поиска.
    \item[Perplexity.] Средняя уверенность в истинных токенах.
    \item[ROUGE.] Пересечение $n$-грамм для суммаризации.
    \item[BLEU.] Совпадение $n$-грамм для перевода.
    \item[Pass@$k$.] Вероятность получить хотя бы одно правильное решение
    среди $k$ попыток.
    \item[Retrieval Recall@$k$.] Доля запросов, для которых нужный
    документ попал в первые $k$ результатов.
    \item[Человеческая оценка.] Полезность, корректность, стиль,
    безопасность и полнота.
\end{description}

Автоматическая метрика должна соответствовать реальной цели. Например,
лексическое совпадение плохо оценивает правильный ответ, сформулированный
другими словами.

\subsection{Формула Pass@$k$}

\begin{proposition}
Пусть получено $n$ решений, из которых $c$ правильных. Если равновероятно
выбрать без возвращения $k$ решений, вероятность того, что среди них
есть хотя бы одно правильное, равна
\[
\operatorname{pass@}k
=
1-
\frac{\binom{n-c}{k}}
{\binom{n}{k}},
\]
где $\binom{n-c}{k}=0$, если $n-c<k$.
\end{proposition}

\begin{proof}
Всего существует
\[
\binom{n}{k}
\]
подмножеств размера $k$. Событие <<нет ни одного правильного решения>>
означает, что все $k$ решений выбраны из $n-c$ неправильных. Таких
подмножеств
\[
\binom{n-c}{k}.
\]
Следовательно,
\[
\Pr(\text{нет правильного})
=
\frac{\binom{n-c}{k}}
{\binom{n}{k}}.
\]
Искомое событие является дополнением:
\[
\Pr(\text{есть правильное})
=
1-
\Pr(\text{нет правильного}),
\]
что и даёт требуемую формулу.
\end{proof}

\subsection{Многокритериальная оценка}

Помимо точности следует измерять:

\begin{itemize}
    \item калибровку уверенности;
    \item устойчивость к переформулировкам;
    \item работу на редких и граничных случаях;
    \item справедливость для разных групп;
    \item токсичность;
    \item фактическую корректность;
    \item отказ от ответа при недостатке данных;
    \item устойчивость к атакующим запросам;
    \item утечку конфиденциальной информации;
    \item эффективность.
\end{itemize}

Высокий средний балл может скрывать критические ошибки в отдельных
категориях.

\subsection{Оценка RAG}

Полезно раздельно измерять:

\begin{enumerate}
    \item находится ли нужный документ;
    \item попадает ли нужный фрагмент в top-$k$;
    \item правильно ли reranker меняет порядок;
    \item поддерживается ли ответ контекстом;
    \item все ли ключевые утверждения имеют источник;
    \item корректно ли модель отказывается при отсутствии данных.
\end{enumerate}

Если оценивать только итоговый ответ, невозможно понять, произошла ошибка
поиска или генерации.

\subsection{Человеческое сравнение}

Два ответа показываются оценщику в случайном порядке. Оценщик выбирает:

\begin{itemize}
    \item ответ A;
    \item ответ B;
    \item ничья;
    \item оба ответа неприемлемы.
\end{itemize}

Критерии должны быть сформулированы заранее. Оценщики должны получать
одинаковые инструкции, а согласованность между ними следует измерять.

LLM-as-a-judge может использоваться для масштабирования оценки, но
судья-модель может иметь:

\begin{itemize}
    \item позиционное смещение;
    \item предпочтение длинных ответов;
    \item предпочтение собственного стиля;
    \item ошибки на узких предметных вопросах;
    \item чувствительность к формулировке критерия.
\end{itemize}

Поэтому автоматического судью калибруют на человеческой выборке.

\subsection{Статистическая надёжность}

Следует сообщать:

\begin{itemize}
    \item число примеров;
    \item среднее значение;
    \item доверительный интервал;
    \item разбивку по категориям;
    \item несколько случайных запусков при стохастической генерации;
    \item параметры декодирования.
\end{itemize}

Для сравнения двух моделей полезна парная оценка на одних и тех же
примерах. Доверительные интервалы можно строить bootstrap-методом:
многократно выбирать примеры с возвращением и пересчитывать метрику.

\subsection{Эксплуатационные метрики}

\begin{description}
    \item[TTFT.] Время до первого токена.
    \item[TPOT.] Среднее время между последующими токенами.
    \item[Latency.] Полное время ответа.
    \item[Throughput.] Запросы или токены в секунду.
    \item[P50, P95, P99.] Квантили задержки.
    \item[Стоимость.] Стоимость запроса или миллиона токенов.
    \item[Utilization.] Использование памяти и вычислительных блоков.
    \item[Error rate.] Доля ошибок, тайм-аутов и нарушений схемы.
\end{description}

% ============================================================
\section{Ограничения и риски}
% ============================================================

\subsection{Галлюцинации}

LLM оптимизируется для вероятного продолжения текста, а не для
доказательства истинности каждого утверждения. Правдоподобная
формулировка может быть фактически неверной.

Причины:

\begin{itemize}
    \item отсутствие нужного знания;
    \item противоречия в обучающих данных;
    \item ошибочное обобщение;
    \item неполный контекст;
    \item стремление выполнить запрос вместо отказа;
    \item ошибки при многошаговом выводе;
    \item неудачное сэмплирование.
\end{itemize}

Меры:

\begin{itemize}
    \item RAG;
    \item инструменты;
    \item цитирование;
    \item проверка утверждений;
    \item ограничение области ответа;
    \item требование сообщать о недостатке данных;
    \item независимые тесты и верификаторы.
\end{itemize}

\subsection{Ограниченность контекста}

Контекстное окно конечно. Даже если документ формально помещается,
модель может:

\begin{itemize}
    \item хуже использовать информацию из середины;
    \item смешивать несколько похожих фактов;
    \item уделять чрезмерное внимание последним инструкциям;
    \item терять связи между удалёнными частями.
\end{itemize}

Увеличение окна не заменяет качественный поиск, структурирование и
сжатие контекста.

\subsection{Запоминание и конфиденциальность}

Большая модель может запоминать редкие и повторяющиеся фрагменты.
Риск увеличивается для:

\begin{itemize}
    \item многочисленных дубликатов;
    \item уникальных идентификаторов;
    \item персональных данных;
    \item секретов и ключей;
    \item небольших специализированных наборов.
\end{itemize}

Меры:

\begin{itemize}
    \item дедупликация;
    \item удаление секретов и персональных данных;
    \item контроль доступа;
    \item privacy-тестирование;
    \item фильтрация выходов;
    \item ограничение логирования;
    \item минимизация передаваемых внешнему сервису данных.
\end{itemize}

\subsection{Смещения}

Модель воспроизводит статистические закономерности данных, включая
стереотипы и неравномерное представление языков и групп.

Средства контроля:

\begin{itemize}
    \item анализ корпуса;
    \item балансировка данных;
    \item специализированное последующее обучение;
    \item многогрупповое тестирование;
    \item человеческий контроль в критических задачах;
    \item документирование известных ограничений.
\end{itemize}

\subsection{Ошибки инструментов и агентов}

Модель может:

\begin{itemize}
    \item выбрать неверный инструмент;
    \item передать неправильные аргументы;
    \item повторить необратимую операцию;
    \item интерпретировать ошибку API как успешный результат;
    \item выполнить инструкцию из недоверенного документа.
\end{itemize}

Защита должна находиться вне LLM:

\begin{itemize}
    \item проверка схем;
    \item allowlist инструментов;
    \item права доступа;
    \item транзакции;
    \item идемпотентность;
    \item подтверждение опасных действий;
    \item песочница;
    \item лимиты времени и числа шагов;
    \item аудит.
\end{itemize}

\subsection{Недетерминированность}

Даже при одинаковых параметрах декодирования возможны небольшие различия
из-за:

\begin{itemize}
    \item случайного сэмплирования;
    \item параллельных вычислений;
    \item порядка редукций;
    \item ограниченной точности;
    \item изменений серверного окружения.
\end{itemize}

Для воспроизводимости фиксируют:

\begin{itemize}
    \item версию модели;
    \item tokenizer;
    \item промпт и диалоговый шаблон;
    \item seed;
    \item температуру, top-$p$ и top-$k$;
    \item stop-последовательности;
    \item версию внешней базы знаний;
    \item версии инструментов.
\end{itemize}

\subsection{Сводная таблица рисков}

\begin{center}
\begin{tabularx}{\textwidth}{>{\bfseries}l X X}
\toprule
Риск & Причина & Основные меры \\
\midrule
Галлюцинации
& Вероятностное продолжение не гарантирует истину
& RAG, инструменты, проверка, отказ при недостатке данных
\\
Утечка данных
& Запоминание и передача чувствительного контекста
& Дедупликация, очистка, контроль доступа, фильтрация
\\
Prompt injection
& Недоверенный текст воспринимается как инструкция
& Изоляция данных, валидация действий, минимум привилегий
\\
Смещения
& Несбалансированный корпус и предпочтения оценщиков
& Аудит данных, групповые тесты, последующее обучение
\\
Ошибки агентов
& Неверные действия и аргументы
& Схемы, подтверждение, sandbox, журналирование
\\
Устаревшие знания
& Параметры фиксируются после обучения
& RAG, поиск, обновляемые базы, continued pretraining
\\
Высокая стоимость
& Большие веса и KV-кэш
& Меньшая модель, GQA, квантование, batching, кэширование
\\
\bottomrule
\end{tabularx}
\end{center}

% ============================================================
\section{Практический конвейер разработки}
% ============================================================

\subsection{Постановка задачи}

До выбора модели определяются:

\begin{enumerate}
    \item входные и выходные данные;
    \item допустимая вероятность ошибки;
    \item цена ложноположительной и ложноотрицательной ошибки;
    \item требования к задержке;
    \item объём запросов;
    \item политика конфиденциальности;
    \item необходимость актуальных знаний;
    \item допустимость человеческой проверки;
    \item критерии успешности.
\end{enumerate}

\subsection{Построение baseline}

Следует начинать с простого решения:

\begin{enumerate}
    \item готовая модель без fine-tuning;
    \item короткий однозначный промпт;
    \item небольшая репрезентативная тестовая выборка;
    \item детерминированное декодирование;
    \item ручной анализ ошибок.
\end{enumerate}

После baseline определяется, что именно ограничивает качество:

\begin{itemize}
    \item отсутствие знаний;
    \item неправильный формат;
    \item недостаточная специализация;
    \item плохой retrieval;
    \item слабая модель;
    \item некорректная постановка задачи.
\end{itemize}

\subsection{Итеративное улучшение}

Рекомендуемый порядок:

\begin{enumerate}
    \item улучшить инструкцию и формат;
    \item добавить few-shot примеры;
    \item добавить программную валидацию;
    \item подключить RAG или инструменты;
    \item улучшить retrieval;
    \item применить SFT или LoRA;
    \item при необходимости выбрать более сильную модель;
    \item оптимизировать инференс;
    \item провести safety- и нагрузочное тестирование.
\end{enumerate}

Fine-tuning не должен быть первым способом исправления проблемы, если
неверные ответы вызваны отсутствием актуальных документов.

\subsection{Развёртывание}

Перед запуском проверяются:

\begin{itemize}
    \item версия модели и токенизатора;
    \item параметры генерации;
    \item права доступа;
    \item политика хранения запросов;
    \item тайм-ауты;
    \item лимиты длины;
    \item обработка ошибок;
    \item fallback-модель или отказ;
    \item мониторинг качества;
    \item возможность отката.
\end{itemize}

\subsection{Мониторинг после запуска}

Необходимо отслеживать:

\begin{itemize}
    \item изменение распределения запросов;
    \item долю отказов;
    \item фактические ошибки;
    \item жалобы пользователей;
    \item задержку и стоимость;
    \item атаки и prompt injection;
    \item устаревание базы знаний;
    \item изменение качества после обновлений.
\end{itemize}

Производственная LLM-система является не только моделью, но и
совокупностью данных, промптов, retrieval, инструментов, политик,
валидаторов и мониторинга.

% ============================================================
\section{Краткая систематизация}
% ============================================================

\subsection{Принципы построения}

\begin{enumerate}
    \item Представить текст токенами.
    \item Разложить вероятность последовательности по цепному правилу.
    \item Использовать Transformer для оценки условных вероятностей.
    \item Ограничить самовнимание каузальной маской.
    \item Обучать модель максимизацией правдоподобия.
    \item Согласовать размер модели, данные и вычислительный бюджет.
    \item Очищать, дедуплицировать и документировать корпус.
    \item Использовать распределённое обучение и смешанную точность.
    \item Выполнить инструктивное и предпочтительное обучение.
    \item Оценивать модель по качеству, безопасности и эффективности.
\end{enumerate}

\subsection{Принципы использования}

\begin{enumerate}
    \item Чётко формулировать задачу и формат.
    \item Не считать уверенный текст доказательством истинности.
    \item Использовать внешние знания для актуальных фактов.
    \item Использовать инструменты для вычислений и действий.
    \item Валидировать структурированный вывод программно.
    \item Ограничивать права агентной системы.
    \item Защищаться от инструкций в недоверенных данных.
    \item Подбирать параметры декодирования под задачу.
    \item Измерять не только среднюю точность, но и критические ошибки.
    \item Постоянно мониторить систему после развёртывания.
\end{enumerate}

\subsection{Основные формулы}

\[
p_\theta(x_{1:n})
=
\prod_{t=1}^{n}p_\theta(x_t\mid x_{<t}).
\]

\[
\mathcal{L}_{\mathrm{NLL}}
=
-\frac{1}{N}
\sum_t
\log p_\theta(x_t\mid x_{<t}).
\]

\[
\operatorname{Attention}(Q,K,V)
=
\softmax
\left(
\frac{QK^\top}{\sqrt{d_k}}+M
\right)V.
\]

\[
\operatorname{PPL}
=
\exp(\mathcal{L}_{\mathrm{NLL}}).
\]

\[
W'=W+\frac{\alpha}{r}BA.
\]

\[
\pi^\ast(y)
=
\frac{
\pi_{\mathrm{ref}}(y)\exp(r(y)/\beta)
}{
\sum_z\pi_{\mathrm{ref}}(z)\exp(r(z)/\beta)
}.
\]

\[
M_{\mathrm{KV}}
\approx
2Lnh_{\mathrm{kv}}d_{\mathrm{head}}b.
\]

\[
N_{\mathrm{params}}
\approx
Vd+
L\left(4d^2+2dd_{\mathrm{ff}}\right).
\]

% ============================================================
\appendix
\section{Краткий словарь терминов}
% ============================================================

\begin{description}
    \item[Checkpoint.] Сохранённое состояние весов и, при необходимости,
    оптимизатора.

    \item[Context window.] Максимальное число токенов, обрабатываемых в
    одном запросе.

    \item[Embedding.] Векторное представление токена или текста.

    \item[Fine-tuning.] Дополнительное обучение предварительно обученной
    модели.

    \item[Hallucination.] Правдоподобный, но неподдержанный или неверный
    вывод модели.

    \item[Inference.] Использование обученной модели без изменения её
    основных весов.

    \item[KV-cache.] Кэш ключей и значений внимания для ранее обработанных
    токенов.

    \item[Logit.] Ненормированная оценка токена перед softmax.

    \item[PEFT.] Parameter-Efficient Fine-Tuning --- адаптация с малым
    числом обучаемых параметров.

    \item[Prompt.] Входная инструкция и контекст модели.

    \item[RAG.] Генерация с извлечением внешних документов.

    \item[Reward model.] Модель, предсказывающая предпочтительность ответа.

    \item[RLHF.] Обучение с подкреплением по обратной связи от человека.

    \item[SFT.] Контролируемое обучение на парах <<запрос--ответ>>.

    \item[Token.] Элемент дискретного словаря токенизатора.

    \item[Tokenizer.] Алгоритм преобразования текста в токены и обратно.

    \item[TTFT.] Время от начала запроса до первого сгенерированного токена.
\end{description}

\end{document}
